\documentclass[11pt]{article}

\usepackage[a4paper,margin=1in]{geometry}
\usepackage{amsmath,amssymb,amsthm,mathtools}
\usepackage{mathrsfs}
\usepackage{hyperref}
\usepackage{enumitem}
\usepackage{microtype}

\hypersetup{
	colorlinks=true,
	linkcolor=blue,
	urlcolor=blue,
	citecolor=blue
}

\title{Lecture Notes: Counting \(k\)-Dimensional Subspaces of an \(n\)-Dimensional Vector Space via Bases}
\author{}
\date{}

\newtheorem{theorem}{Theorem}[section]
\newtheorem{proposition}[theorem]{Proposition}
\newtheorem{lemma}[theorem]{Lemma}
\newtheorem{corollary}[theorem]{Corollary}
\theoremstyle{definition}
\newtheorem{definition}[theorem]{Definition}
\newtheorem{example}[theorem]{Example}
\newtheorem{exercise}[theorem]{Exercise}
\theoremstyle{remark}
\newtheorem{remark}[theorem]{Remark}

\newcommand{\F}{\mathbf{F}}
\newcommand{\R}{\mathbf{R}}
\newcommand{\C}{\mathbf{C}}
\newcommand{\Q}{\mathbf{Q}}
\newcommand{\Span}{\operatorname{span}}
\newcommand{\GL}{\mathrm{GL}}
\newcommand{\Gr}{\mathrm{Gr}}
\newcommand{\qbinom}[2]{\genfrac{[}{]}{0pt}{}{#1}{#2}_q}

\begin{document}
	
	\maketitle
	
	\tableofcontents
	
	\section{Introduction}
	
	Let \(V\) be an \(n\)-dimensional vector space over a field \(K\). We want to answer the following question:
	
	\begin{quote}
		How many distinct \(k\)-dimensional subspaces does \(V\) have?
	\end{quote}
	
	The answer depends on the field \(K\).
	
	\begin{itemize}
		\item If \(K=\F_q\) is a finite field with \(q\) elements, then the number of \(k\)-dimensional subspaces is finite and is given by the \emph{Gaussian binomial coefficient}.
		\item If \(K\) is infinite, then for \(0<k<n\) there are infinitely many \(k\)-dimensional subspaces.
	\end{itemize}
	
	In these notes, we develop the finite-field formula using the most concrete approach: \emph{counting bases}. The key idea is simple:
	
	\begin{quote}
		Count all ordered bases of all \(k\)-dimensional subspaces of \(\F_q^n\), and then divide by the number of ordered bases of one fixed \(k\)-dimensional subspace.
	\end{quote}
	
	This basis-counting method is the most direct way to understand why the formula has the shape it does.
	
	\section{Why We May Work in \(\F_q^n\)}
	
	Suppose \(V\) is an \(n\)-dimensional vector space over \(\F_q\). After choosing a basis of \(V\), we obtain an isomorphism
	\[
	V \cong \F_q^n.
	\]
	Since vector space isomorphisms preserve subspaces and dimensions, counting \(k\)-dimensional subspaces of \(V\) is equivalent to counting \(k\)-dimensional subspaces of \(\F_q^n\).
	
	\begin{proposition}
		If \(T:V\to W\) is a vector space isomorphism, then the map
		\[
		U\longmapsto T(U)
		\]
		gives a bijection between the \(k\)-dimensional subspaces of \(V\) and the \(k\)-dimensional subspaces of \(W\).
	\end{proposition}
	
	\begin{proof}
		If \(U\subseteq V\) is a subspace, then \(T(U)\subseteq W\) is a subspace. Since \(T|_U:U\to T(U)\) is an isomorphism, we have
		\[
		\dim T(U)=\dim U.
		\]
		The inverse correspondence is given by \(X\mapsto T^{-1}(X)\), so the map is bijective.
	\end{proof}
	
	\begin{remark}
		Thus the number of \(k\)-dimensional subspaces depends only on \(n\) and the field, not on the specific vector space.
	\end{remark}
	
	So from now on we assume
	\[
	V=\F_q^n.
	\]
	
	\section{Ordinary Binomial Coefficients as Motivation}
	
	Before counting subspaces, recall a familiar counting problem.
	
	\begin{quote}
		How many \(k\)-element subsets does an \(n\)-element set have?
	\end{quote}
	
	The answer is
	\[
	\binom{n}{k}.
	\]
	
	The finite-field analogue replaces subsets by subspaces:
	\[
	\text{\(k\)-element subsets of an \(n\)-element set}
	\quad \longleftrightarrow \quad
	\text{\(k\)-dimensional subspaces of \(\F_q^n\)}.
	\]
	
	This leads to the \emph{Gaussian binomial coefficient}, also called the \emph{\(q\)-binomial coefficient}.
	
	\section{Statement of the Main Formula}
	
	\begin{definition}
		For integers \(0\le k\le n\), the \emph{Gaussian binomial coefficient} is
		\[
		\qbinom{n}{k}
		=
		\frac{(q^n-1)(q^n-q)(q^n-q^2)\cdots(q^n-q^{k-1})}
		{(q^k-1)(q^k-q)(q^k-q^2)\cdots(q^k-q^{k-1})}.
		\]
		Equivalently,
		\[
		\qbinom{n}{k}
		=
		\prod_{i=0}^{k-1}\frac{q^n-q^i}{q^k-q^i}.
		\]
	\end{definition}
	
	\begin{theorem}[Main Theorem]
		Let \(V\) be an \(n\)-dimensional vector space over \(\F_q\). Then the number of distinct \(k\)-dimensional subspaces of \(V\) is
		\[
		\boxed{
			\qbinom{n}{k}
			=
			\prod_{i=0}^{k-1}\frac{q^n-q^i}{q^k-q^i}
		}.
		\]
	\end{theorem}
	
	Our goal is to prove this theorem using bases.
	
	\section{The Basis-Counting Strategy}
	
	To count \(k\)-dimensional subspaces of \(\F_q^n\), we proceed in two stages.
	
	\begin{enumerate}[label=\arabic*.]
		\item Count all ordered linearly independent \(k\)-tuples
		\[
		(v_1,\dots,v_k)
		\]
		in \(\F_q^n\). Each such tuple is an ordered basis of some \(k\)-dimensional subspace.
		\item Determine how many ordered bases a single fixed \(k\)-dimensional subspace has.
	\end{enumerate}
	
	Then:
	\[
	\#\{\text{\(k\)-dimensional subspaces}\}
	=
	\frac{\#\{\text{ordered bases of all \(k\)-subspaces}\}}
	{\#\{\text{ordered bases of one \(k\)-subspace}\}}.
	\]
	
	This is the central idea of the whole argument.
	
	\section{Counting Ordered Independent \(k\)-Tuples}
	
	We now count the ordered \(k\)-tuples
	\[
	(v_1,\dots,v_k)\in (\F_q^n)^k
	\]
	that are linearly independent.
	
	\begin{lemma}
		The number of ordered linearly independent \(k\)-tuples in \(\F_q^n\) is
		\[
		(q^n-1)(q^n-q)(q^n-q^2)\cdots(q^n-q^{k-1}).
		\]
	\end{lemma}
	
	\begin{proof}
		We choose the vectors one at a time.
		
		\begin{itemize}
			\item The first vector \(v_1\) can be any nonzero vector of \(\F_q^n\). Since \(\F_q^n\) has \(q^n\) vectors total, there are
			\[
			q^n-1
			\]
			choices for \(v_1\).
			
			\item The second vector \(v_2\) must not lie in \(\Span(v_1)\). Since \(\Span(v_1)\) is a 1-dimensional subspace, it has exactly \(q\) elements. Therefore there are
			\[
			q^n-q
			\]
			choices for \(v_2\).
			
			\item The third vector \(v_3\) must not lie in \(\Span(v_1,v_2)\). Since \(v_1,v_2\) are linearly independent, this span has dimension \(2\), hence \(q^2\) elements. So there are
			\[
			q^n-q^2
			\]
			choices for \(v_3\).
		\end{itemize}
		
		Continuing this way, when choosing \(v_i\), the previously chosen vectors span a subspace of dimension \(i-1\), which has \(q^{i-1}\) elements. Thus the number of choices for \(v_i\) is
		\[
		q^n-q^{i-1}.
		\]
		
		Multiplying over \(i=1,\dots,k\), we obtain
		\[
		(q^n-1)(q^n-q)(q^n-q^2)\cdots(q^n-q^{k-1}).
		\]
	\end{proof}
	
	\begin{remark}
		This count is sometimes called the number of ordered \(k\)-frames in \(\F_q^n\).
	\end{remark}
	
	\section{Counting Ordered Bases of a Fixed \(k\)-Dimensional Subspace}
	
	Now fix one \(k\)-dimensional subspace \(W\subseteq \F_q^n\). We ask: how many ordered bases does \(W\) have?
	
	\begin{lemma}
		A fixed \(k\)-dimensional subspace \(W\) has exactly
		\[
		(q^k-1)(q^k-q)(q^k-q^2)\cdots(q^k-q^{k-1})
		\]
		ordered bases.
	\end{lemma}
	
	\begin{proof}
		Since \(W\) is \(k\)-dimensional, it is isomorphic to \(\F_q^k\). So counting ordered bases of \(W\) is the same as counting ordered linearly independent \(k\)-tuples in \(\F_q^k\).
		
		Applying the previous lemma with \(n=k\), we get
		\[
		(q^k-1)(q^k-q)(q^k-q^2)\cdots(q^k-q^{k-1}).
		\]
	\end{proof}
	
	\begin{remark}
		This number is also the order of the general linear group \(\GL_k(\F_q)\), since an invertible \(k\times k\) matrix is exactly an ordered basis of \(\F_q^k\).
	\end{remark}
	
	\section{Proof of the Main Theorem}
	
	We now combine the two counts.
	
	\begin{proof}[Proof of the Main Theorem]
		Every ordered linearly independent \(k\)-tuple in \(\F_q^n\) spans a unique \(k\)-dimensional subspace. Conversely, each \(k\)-dimensional subspace contributes exactly as many ordered linearly independent \(k\)-tuples as it has ordered bases, namely
		\[
		(q^k-1)(q^k-q)(q^k-q^2)\cdots(q^k-q^{k-1}).
		\]
		
		Therefore,
		\[
		\#\{\text{\(k\)-dimensional subspaces of }\F_q^n\}
		=
		\frac{(q^n-1)(q^n-q)\cdots(q^n-q^{k-1})}
		{(q^k-1)(q^k-q)\cdots(q^k-q^{k-1})}.
		\]
		
		By definition, this is exactly
		\[
		\qbinom{n}{k}.
		\]
	\end{proof}
	
	\section{Why the Formula Has This Shape}
	
	It is worth pausing to interpret the formula:
	\[
	\frac{(q^n-1)(q^n-q)\cdots(q^n-q^{k-1})}
	{(q^k-1)(q^k-q)\cdots(q^k-q^{k-1})}.
	\]
	
	\begin{itemize}
		\item The numerator counts all ordered bases of all \(k\)-dimensional subspaces inside \(\F_q^n\).
		\item The denominator counts how many ordered bases correspond to one fixed \(k\)-dimensional subspace.
	\end{itemize}
	
	So the quotient removes the overcounting caused by the fact that the same subspace has many different ordered bases.
	
	This is exactly analogous to ordinary counting arguments where one counts objects with labels and then divides by the number of labelings.
	
	\section{Examples}
	
	\begin{example}[Lines in \(\F_q^3\)]
		A \(1\)-dimensional subspace is a line through the origin. The number of such lines is
		\[
		\qbinom{3}{1}
		=
		\frac{q^3-1}{q-1}
		=
		q^2+q+1.
		\]
	\end{example}
	
	\begin{example}[Planes in \(\F_q^3\)]
		A \(2\)-dimensional subspace is a plane through the origin. The number of such planes is
		\[
		\qbinom{3}{2}
		=
		\qbinom{3}{1}
		=
		q^2+q+1.
		\]
	\end{example}
	
	\begin{example}[All lines in \(\F_2^2\)]
		The space \(\F_2^2\) has vectors
		\[
		(0,0),\ (1,0),\ (0,1),\ (1,1).
		\]
		Its nonzero vectors are
		\[
		(1,0),\ (0,1),\ (1,1).
		\]
		Each 1-dimensional subspace contains exactly one nonzero vector, so there are exactly three lines:
		\[
		\Span(1,0),\qquad
		\Span(0,1),\qquad
		\Span(1,1).
		\]
		This agrees with
		\[
		\qbinom{2}{1}
		=
		\frac{2^2-1}{2-1}
		=
		3.
		\]
	\end{example}
	
	\begin{example}[Two-dimensional subspaces of \(\F_2^4\)]
		Take \(q=2\), \(n=4\), \(k=2\). Then
		\[
		\qbinom{4}{2}
		=
		\frac{(2^4-1)(2^4-2)}{(2^2-1)(2^2-2)}
		=
		\frac{15\cdot 14}{3\cdot 2}
		=
		35.
		\]
		So \(\F_2^4\) has exactly \(35\) distinct 2-dimensional subspaces.
	\end{example}
	
	\section{A Fully Concrete Example: \(\F_2^3\)}
	
	Let us count the 2-dimensional subspaces of \(\F_2^3\) directly by bases.
	
	The formula predicts:
	\[
	\qbinom{3}{2}\Big|_{q=2}
	=
	\frac{(2^3-1)(2^3-2)}{(2^2-1)(2^2-2)}
	=
	\frac{7\cdot 6}{3\cdot 2}
	=
	7.
	\]
	
	Now let us interpret this.
	
	\subsection*{Step 1: Count ordered bases of 2-dimensional subspaces of \(\F_2^3\)}
	
	To choose an ordered basis \((v_1,v_2)\):
	
	\begin{itemize}
		\item \(v_1\) can be any nonzero vector: \(7\) choices.
		\item \(v_2\) must not lie in \(\Span(v_1)\). Over \(\F_2\), the span of a nonzero vector contains only \(0\) and the vector itself, so there are \(6\) choices.
	\end{itemize}
	
	Thus there are
	\[
	7\cdot 6=42
	\]
	ordered bases of 2-dimensional subspaces in \(\F_2^3\).
	
	\subsection*{Step 2: Count ordered bases of one fixed 2-dimensional subspace}
	
	A 2-dimensional subspace is isomorphic to \(\F_2^2\). In \(\F_2^2\),
	
	\begin{itemize}
		\item the first basis vector has \(3\) choices,
		\item the second has \(2\) choices.
	\end{itemize}
	
	So one 2-dimensional subspace has
	\[
	3\cdot 2=6
	\]
	ordered bases.
	
	\subsection*{Step 3: Divide}
	
	Therefore the number of distinct 2-dimensional subspaces is
	\[
	\frac{42}{6}=7.
	\]
	
	\section{The General Linear Group and Bases}
	
	The basis-counting argument is closely related to the size of the general linear group.
	
	\begin{proposition}
		The order of \(\GL_k(\F_q)\) is
		\[
		|\GL_k(\F_q)|
		=
		(q^k-1)(q^k-q)(q^k-q^2)\cdots(q^k-q^{k-1}).
		\]
	\end{proposition}
	
	\begin{proof}
		An invertible \(k\times k\) matrix is exactly an ordered basis of \(\F_q^k\), written as its columns. Counting such ordered bases gives the formula.
	\end{proof}
	
	\begin{remark}
		So the denominator in the Gaussian coefficient is exactly \(|\GL_k(\F_q)|\). This means:
		\[
		\qbinom{n}{k}
		=
		\frac{\text{number of ordered \(k\)-frames in \(\F_q^n\)}}{|\GL_k(\F_q)|}.
		\]
	\end{remark}
	
	\section{Special Cases}
	
	\begin{proposition}
		For every \(n\) and \(q\),
		\begin{align*}
			\qbinom{n}{0} &= 1,\\
			\qbinom{n}{1} &= \frac{q^n-1}{q-1}=1+q+q^2+\cdots+q^{n-1},\\
			\qbinom{n}{n-1} &= \qbinom{n}{1},\\
			\qbinom{n}{n} &= 1.
		\end{align*}
	\end{proposition}
	
	\begin{proof}
		The cases \(k=0\) and \(k=n\) are obvious. For \(k=1\), count nonzero vectors and divide by the number \(q-1\) of nonzero scalar multiples of a given vector. The case \(k=n-1\) follows from symmetry.
	\end{proof}
	
	\section{Symmetry}
	
	\begin{proposition}
		For \(0\le k\le n\),
		\[
		\qbinom{n}{k}=\qbinom{n}{n-k}.
		\]
	\end{proposition}
	
	\begin{proof}
		This follows directly from the product formulas. It also reflects a geometric duality between \(k\)-dimensional and \((n-k)\)-dimensional objects.
	\end{proof}
	
	\section{Another Form of the Gaussian Binomial Coefficient}
	
	The Gaussian coefficient may also be written as
	\[
	\qbinom{n}{k}
	=
	\frac{(q^n-1)(q^{n-1}-1)\cdots(q^{n-k+1}-1)}
	{(q^k-1)(q^{k-1}-1)\cdots(q-1)}.
	\]
	
	This form is often convenient in algebraic manipulations and makes the symmetry
	\[
	\qbinom{n}{k}=\qbinom{n}{n-k}
	\]
	more transparent.
	
	\section{The \(q\to 1\) Philosophy}
	
	The Gaussian binomial coefficient is called a \(q\)-analogue of the ordinary binomial coefficient because, after simplification as a polynomial in \(q\), setting \(q=1\) gives the ordinary binomial coefficient.
	
	\begin{example}
		Since
		\[
		\qbinom{3}{1}=q^2+q+1,
		\]
		substituting \(q=1\) gives
		\[
		1+1+1=3=\binom{3}{1}.
		\]
	\end{example}
	
	So subsets are the \(q=1\) shadow of subspaces over finite fields.
	
	\section{Infinite Fields}
	
	What happens if the field is infinite?
	
	\begin{proposition}
		If \(K\) is an infinite field and \(V\) is an \(n\)-dimensional vector space over \(K\), then for \(0<k<n\), the number of \(k\)-dimensional subspaces of \(V\) is infinite.
	\end{proposition}
	
	\begin{proof}
		It is enough to consider \(k=1\). In \(K^2\), every line through the origin is a 1-dimensional subspace. Since \(K\) is infinite, there are infinitely many such lines. Therefore there are infinitely many \(k\)-dimensional subspaces in general for \(0<k<n\).
	\end{proof}
	
	\begin{remark}
		So the finite formula \(\qbinom{n}{k}\) is special to finite fields.
	\end{remark}
	
	\section{Grassmannians}
	
	\begin{definition}
		The set of all \(k\)-dimensional subspaces of an \(n\)-dimensional vector space is called the \emph{Grassmannian}, denoted
		\[
		\Gr(k,n).
		\]
	\end{definition}
	
	Over \(\F_q\), the Grassmannian is a finite set, and
	\[
	|\Gr(k,n)|=\qbinom{n}{k}.
	\]
	
	So Gaussian binomial coefficients count the points of Grassmannians over finite fields.
	
	\section{Conceptual Summary}
	
	The basis-counting proof rests on the following facts:
	
	\begin{enumerate}[label=\arabic*.]
		\item A \(k\)-dimensional subspace is determined by any ordered basis of it.
		\item The number of ordered linearly independent \(k\)-tuples in \(\F_q^n\) is
		\[
		(q^n-1)(q^n-q)\cdots(q^n-q^{k-1}).
		\]
		\item The number of ordered bases of one fixed \(k\)-dimensional space is
		\[
		(q^k-1)(q^k-q)\cdots(q^k-q^{k-1}).
		\]
		\item Dividing gives the number of distinct \(k\)-dimensional subspaces.
	\end{enumerate}
	
	Thus
	\[
	\boxed{
		\#\{\text{\(k\)-dimensional subspaces of }\F_q^n\}
		=
		\frac{(q^n-1)(q^n-q)\cdots(q^n-q^{k-1})}
		{(q^k-1)(q^k-q)\cdots(q^k-q^{k-1})}
		=
		\qbinom{n}{k}.
	}
	\]
	
	\section{Final Theorem}
	
	\begin{theorem}[Final Form]
		Let \(V\) be an \(n\)-dimensional vector space over \(\F_q\). Then the number of distinct \(k\)-dimensional subspaces of \(V\) is
		\[
		\boxed{
			\qbinom{n}{k}
			=
			\prod_{i=0}^{k-1}\frac{q^n-q^i}{q^k-q^i}
		}.
		\]
		This formula is obtained by counting ordered bases of \(k\)-dimensional subspaces and dividing by the number of ordered bases of one fixed \(k\)-dimensional subspace.
	\end{theorem}
	
	\section{Exercises}
	
	\begin{exercise}
		Prove directly that an \(r\)-dimensional subspace of \(\F_q^n\) has exactly \(q^r\) elements.
	\end{exercise}
	
	\begin{exercise}
		Show that the number of ordered bases of \(\F_q^k\) is \(|\GL_k(\F_q)|\).
	\end{exercise}
	
	\begin{exercise}
		Use basis-counting to compute the number of lines in \(\F_q^4\).
	\end{exercise}
	
	\begin{exercise}
		Compute \(\qbinom{4}{2}\) when \(q=2\).
	\end{exercise}
	
	\begin{exercise}
		Compute \(\qbinom{4}{2}\) when \(q=3\).
	\end{exercise}
	
	\begin{exercise}
		Show that
		\[
		\qbinom{n}{1}=1+q+q^2+\cdots+q^{n-1}.
		\]
	\end{exercise}
	
	\begin{exercise}
		Prove that
		\[
		\qbinom{n}{k}=\qbinom{n}{n-k}.
		\]
	\end{exercise}
	
	\begin{exercise}
		Explain why the basis-counting proof fails over an infinite field if one tries to get a finite number.
	\end{exercise}
	
	\section{Selected Answers}
	
	\begin{itemize}
		\item For \(q=2\),
		\[
		\qbinom{4}{2}
		=
		\frac{(2^4-1)(2^4-2)}{(2^2-1)(2^2-2)}
		=
		\frac{15\cdot 14}{3\cdot 2}
		=
		35.
		\]
		
		\item For \(q=3\),
		\[
		\qbinom{4}{2}
		=
		\frac{(3^4-1)(3^4-3)}{(3^2-1)(3^2-3)}
		=
		\frac{80\cdot 78}{8\cdot 6}
		=
		130.
		\]
		
		\item The number of lines in \(\F_q^4\) is
		\[
		\qbinom{4}{1}
		=
		\frac{q^4-1}{q-1}
		=
		1+q+q^2+q^3.
		\]
	\end{itemize}
	
	\section*{One-Sentence Takeaway}
	
	Over a finite field \(\F_q\), the number of distinct \(k\)-dimensional subspaces of an \(n\)-dimensional vector space is the Gaussian binomial coefficient \(\qbinom{n}{k}\), and the most direct proof comes from counting ordered bases.
	
\end{document}